% =============================================================================
% CAPÍTULO 3 — GUÍA OPERATIVA: PERFIL PROFESIONAL  (generado por tools/gen_ch3.py)
% =============================================================================
\chapter{Guía Operativa: Perfil Profesional}
\label{ch:profesional}
\resumenCapitulo{Este capítulo guía a la persona usuaria con perfil Profesional en la
construcción y gestión de su portafolio digital: proyectos, experiencia laboral,
formación académica, habilidades y currículum, así como en la configuración y la
privacidad de su perfil público y el acceso desde dispositivos móviles.}

\section{Visión General del Espacio de Trabajo (Dashboard)}
\label{sec:prof-dashboard}
Tras iniciar sesión como Profesional, accederá a su espacio de trabajo. Desde el
\textbf{panel lateral izquierdo} se navega a todas las secciones del portafolio:
\campo{Mi portafolio}, \campo{Habilidades}, \campo{Proyectos}, \campo{Experiencia},
\campo{Educación}, \campo{CV Studio}, \campo{Conexiones}, \campo{Chat} y
\campo{Configuración}. El área central muestra el contenido de la sección activa y los
indicadores resumen (número de proyectos, habilidades destacadas, etc.).

\begin{AvisoNota}
El portafolio se construye por bloques independientes. Puede completarlos en el orden
que prefiera; cada bloque se guarda por separado y se refleja de inmediato en su perfil
público.
\end{AvisoNota}

\section{Gestión Integral del Portafolio}
\label{sec:prof-portafolio}
El portafolio reúne sus proyectos, su experiencia, su formación y sus habilidades. En
esta sección se documenta la gestión completa de cada bloque, paso a paso.

\subsection{Creación de un Nuevo Proyecto}
\label{subsec:prof-crear-proyecto}
La creación de un proyecto se realiza mediante un asistente de tres pasos: General,
Técnico y Media.

\subsubsection{Paso 1: Navegación al formulario de creación}

\captura{assets/img/ethoshub-profesional/03_proyectos/CrearProyecto1.png}{Al acceder a la sección «Proyectos» desde el panel lateral izquierdo, observamos la vista de gestión principal con las opciones de filtrado y el botón para iniciar un nuevo proyecto.}
\captura{assets/img/ethoshub-profesional/03_proyectos/CrearProyecto2.png}{Al visualizar el área central de la vista «Proyectos y Evidencias», identificamos el botón principal «Crear nuevo proyecto» que permite iniciar el proceso de registro.}

\subsubsection{Paso 2: Llenado de campos obligatorios (General y Técnico)}

\captura{assets/img/ethoshub-profesional/03_proyectos/CrearProyecto3.png}{Al abrir el formulario de creación, se nos presenta el paso 1, denominado «General», donde se solicita el título del proyecto (en este caso, «Nexus Auth»).}
\captura{assets/img/ethoshub-profesional/03_proyectos/CrearProyecto4.png}{Al observar el campo de categoría dentro del paso «General», podemos seleccionar el tipo de proyecto, siendo «Web» la opción elegida para este registro.}
\captura{assets/img/ethoshub-profesional/03_proyectos/CrearProyecto5.png}{Al completar el campo de descripción, detallamos el propósito del proyecto («Sistema centralizado de autenticación OAuth2»).}
\captura{assets/img/ethoshub-profesional/03_proyectos/CrearProyecto6.png}{Confirma la correcta entrada de la descripción del proyecto en el formulario, manteniendo la coherencia con los datos solicitados.}
\captura{assets/img/ethoshub-profesional/03_proyectos/CrearProyecto7.png}{Al interactuar con el campo «Estado Inicial», se despliega un menú desplegable para definir la fase actual del proyecto.}
\captura{assets/img/ethoshub-profesional/03_proyectos/CrearProyecto8.png}{Muestra las opciones disponibles dentro del selector de estado («Borrador», «En progreso», «Completado», «Archivado»), permitiendo definir el estatus inicial correctamente.}
\captura{assets/img/ethoshub-profesional/03_proyectos/CrearProyecto9.png}{Al definir el rol del usuario en el proyecto, el campo «Rol en el proyecto» permite especificar la responsabilidad, la cual se ha configurado como «Lead Dev».}
\captura{assets/img/ethoshub-profesional/03_proyectos/CrearProyecto10.png}{Al configurar la visibilidad, la opción «Proyecto público» cuenta con un interruptor (toggle) que, al activarse, permite que el proyecto sea visible para otros usuarios.}
\captura{assets/img/ethoshub-profesional/03_proyectos/CrearProyecto11.png}{Al configurar la visibilidad, identificamos el interruptor «Destacado en portafolio», que permite determinar si el proyecto se mostrará en la sección principal.}
\captura{assets/img/ethoshub-profesional/03_proyectos/CrearProyecto12.png}{Al interactuar con el formulario, hacemos clic en el botón «Continuar» para avanzar desde la pantalla de información general hacia los detalles técnicos.}
\captura{assets/img/ethoshub-profesional/03_proyectos/CrearProyecto13.png}{Al observar el formulario de información general, identificamos la opción «Cancelar» para interrumpir el proceso de creación.}
\captura{assets/img/ethoshub-profesional/03_proyectos/CrearProyecto14.png}{Al acceder al paso 2 del asistente, se nos presenta la pantalla «Técnico», donde se solicitan los detalles de fechas, repositorio y tecnologías.}
\captura{assets/img/ethoshub-profesional/03_proyectos/CrearProyecto15.png}{Al interactuar con el campo «Fecha de inicio», se activa el selector de calendario para configurar el inicio del proyecto.}
\captura{assets/img/ethoshub-profesional/03_proyectos/CrearProyecto16.png}{Confirma la selección del «02/05/2024» como la fecha de inicio establecida en el formulario.}
\captura{assets/img/ethoshub-profesional/03_proyectos/CrearProyecto17.png}{Al seleccionar el campo «Fecha de fin», se despliega la interfaz del calendario para definir la conclusión del proyecto.}
\captura{assets/img/ethoshub-profesional/03_proyectos/CrearProyecto18.png}{Ilustra la fecha de fin seleccionada como «02/03/2025» mediante el uso del calendario.}
\captura{assets/img/ethoshub-profesional/03_proyectos/CrearProyecto19.png}{Al observar el campo «Repositorio GitHub», confirmamos que la URL del repositorio ha sido ingresada correctamente.}
\captura{assets/img/ethoshub-profesional/03_proyectos/CrearProyecto20.png}{Al visualizar el área de «Tecnologías utilizadas» en el paso 2, observamos la inclusión de añadir tecnologias como ser «React» y el botón para añadir tecnologías adicionales.}
\captura{assets/img/ethoshub-profesional/03_proyectos/CrearProyecto21.png}{Observamos el campo de texto «Resultados obtenidos» donde se detalla el impacto del proyecto.}
\captura{assets/img/ethoshub-profesional/03_proyectos/CrearProyecto22.png}{Al interactuar con los controles de navegación, el botón «Atrás» permite regresar a la configuración anterior.}
\captura{assets/img/ethoshub-profesional/03_proyectos/CrearProyecto23.png}{Tambien podemos observar el botón «Cancelar» que permite finalizar la edición en el paso 2 y cerrar el formulario.}
\captura{assets/img/ethoshub-profesional/03_proyectos/CrearProyecto24.png}{Tambien esta el boton «Continuar» que nos permite avanzar a la siguiente seccion del formulario.}

\subsubsection{Paso 3: Carga de recursos multimedia y confirmación}

\captura{assets/img/ethoshub-profesional/03_proyectos/CrearProyecto25.png}{Al iniciar el paso 3 «Media», se nos presenta la interfaz para elegir una portada predefinida o cargar archivos.}
\captura{assets/img/ethoshub-profesional/03_proyectos/CrearProyecto26.png}{Al explorar las portadas predefinidas, seleccionamos la opción «Violet Space» para personalizar el aspecto del proyecto.}
\captura{assets/img/ethoshub-profesional/03_proyectos/CrearProyecto27.png}{Al enfocarnos en la sección de «Imagen personalizada», identificamos el área habilitada para subir o arrastrar archivos.}
\captura{assets/img/ethoshub-profesional/03_proyectos/CrearProyecto28.png}{Confirma la carga exitosa de una imagen personalizada, mostrando el logo de «NEXUS» dentro del formulario.}
\captura{assets/img/ethoshub-profesional/03_proyectos/CrearProyecto29.png}{Al utilizar la opción de carga por enlace, el campo «URL de imagen» permite integrar recursos externos.}
\captura{assets/img/ethoshub-profesional/03_proyectos/CrearProyecto30.png}{Ilustra el resultado al presionar «Previsualizar», mostrando la imagen cargada a través de la URL proporcionada.}
\captura{assets/img/ethoshub-profesional/03_proyectos/CrearProyecto31.png}{Tambien observamos el campo «Enlaces y Recursos» que permite integrar plataformas externas como YouTube o Figma mediante el botón «Añadir».}
\captura{assets/img/ethoshub-profesional/03_proyectos/CrearProyecto32.png}{Muestra la correcta incorporación de un enlace a un video de YouTube, incluyendo la opción de eliminar el recurso mediante el icono de la papelera.}
\captura{assets/img/ethoshub-profesional/03_proyectos/CrearProyecto33.png}{Al descender en el formulario, identificamos el área de «Documentos», donde el sistema habilita la carga de archivos PDF mediante arrastrar y soltar.}
\captura{assets/img/ethoshub-profesional/03_proyectos/CrearProyecto34.png}{Al utilizar la función para documentos externos, el campo específico permite vincular archivos de Google Drive mediante el botón «Añadir».}
\captura{assets/img/ethoshub-profesional/03_proyectos/CrearProyecto35.png}{Confirma la correcta gestión de múltiples documentos, mostrando tanto el archivo cargado localmente como el enlace externo bajo la categoría de documentos.}
\captura{assets/img/ethoshub-profesional/03_proyectos/CrearProyecto36.png}{Al finalizar la configuración en el paso 3, el botón «Crear proyecto» queda habilitado para guardar toda la información ingresada.}
\captura{assets/img/ethoshub-profesional/03_proyectos/CrearProyecto37.png}{Muestra la confirmación visual en la parte superior derecha, indicando que el proyecto se ha registrado exitosamente en el sistema.}

\subsection{Edición y Actualización de Proyectos Existentes}
\label{subsec:prof-editar-proyecto}
Una vez creado, el proyecto aparece en el listado de \campo{Proyectos y Evidencias},
donde puede consultarse, filtrarse y editarse.

\captura{assets/img/ethoshub-profesional/03_proyectos/CrearProyecto38.png}{Al volver a la vista principal, observamos el listado actualizado donde aparece el proyecto «Nexus Auth» junto a los controles de filtrado.}
\captura{assets/img/ethoshub-profesional/03_proyectos/CrearProyecto39.png}{Al observar la vista general de «Proyectos y Evidencias», identificamos algunos filtros para clasificar proyectos.}
\captura{assets/img/ethoshub-profesional/03_proyectos/CrearProyecto40.png}{Tambien al observar la vista general de «Proyectos y Evidencias», identificamos el filtro desplegable para clasificar los proyectos según su estado actual.}
\captura{assets/img/ethoshub-profesional/03_proyectos/CrearProyecto41.png}{Al observar la vista general de «Proyectos y Evidencias», identificamos el filtro de «Destacados primero».}
\captura{assets/img/ethoshub-profesional/03_proyectos/CrearProyecto42.png}{Al interactuar con el selector de ordenamiento en la parte superior derecha, identificamos el campo «Destacados primero» que permite organizar la visualización.}
\captura{assets/img/ethoshub-profesional/03_proyectos/CrearProyecto43.png}{Al observar el área de estadísticas del panel, identificamos el recuento de proyectos registrados en la plataforma.}
\captura{assets/img/ethoshub-profesional/03_proyectos/CrearProyecto44.png}{Al ubicar el botón principal de la sección «Proyectos y Evidencias», confirmamos el acceso para crear un nuevo proyecto mediante «Agregar proyecto».}
\captura{assets/img/ethoshub-profesional/03_proyectos/CrearProyecto45.png}{Al iniciar el formulario de creacion, se nos presenta el asistente de registro con el primer paso configurado en «General».}
\captura{assets/img/ethoshub-profesional/03_proyectos/CrearProyecto46.png}{Se puede observar que en el proyecto creado existe 3 funcionalidades, la primera y la que se ve es poder «Ver a detalle» un proyecto creado.}
\captura{assets/img/ethoshub-profesional/03_proyectos/CrearProyecto47.png}{Ilustra la vista detallada del proyecto, presentando la información técnica, los resultados, el impacto y la gestión de documentos asociados.}
\captura{assets/img/ethoshub-profesional/03_proyectos/CrearProyecto48.png}{Al volver a la vista principal, identificamos el icono del lápiz sobre la tarjeta del proyecto, el cual permite acceder a la función de edición.}
\captura{assets/img/ethoshub-profesional/03_proyectos/CrearProyecto49.png}{Al activar la edición, el formulario cambia su título a «Editar proyecto» y permite modificar los datos que deseemos antes de finalizar con «Actualizar proyecto».}

\subsection{Eliminación o Depuración de Elementos del Portafolio}
\label{subsec:prof-eliminar-proyecto}

\captura{assets/img/ethoshub-profesional/03_proyectos/CrearProyecto50.png}{Al volver a la vista principal, identificamos el icono del «basurero» que nos permite eliminar un proyecto.}
\captura{assets/img/ethoshub-profesional/03_proyectos/CrearProyecto51.png}{Al confirmar la eliminación de un proyecto, aparece un diálogo de advertencia que solicita al usuario verificar la acción antes de proceder.}
\captura{assets/img/ethoshub-profesional/03_proyectos/CrearProyecto52.png}{Una vez completada la eliminación, el sistema muestra una notificación de confirmación en la esquina superior derecha e indica que actualmente no existen proyectos registrados.}

\begin{AvisoAdvertencia}
La eliminación de un proyecto es permanente. Asegúrese de no necesitar su información
antes de confirmar el borrado.
\end{AvisoAdvertencia}

\subsection{Registro de Experiencia Laboral}
\label{subsec:prof-experiencia}
El bloque \campo{Experiencia} permite documentar cada puesto de trabajo con sus fechas,
ubicación, tecnologías y evidencias visuales.

\captura{assets/img/ethoshub-profesional/04_experiencia/profesional_experiencia_laboral_1.png}{Para ir al apartado de experiencia laboral debe hacer clic al boton «Experiencia».}
\captura{assets/img/ethoshub-profesional/04_experiencia/profesional_experiencia_laboral_2.png}{Despues de darle al boton aparece la pantalla.}
\captura{assets/img/ethoshub-profesional/04_experiencia/profesional_experiencia_laboral_3.png}{Para poder registrar una experiencia debe darle al boton de «añadir».}
\captura{assets/img/ethoshub-profesional/04_experiencia/profesional_experiencia_laboral_4.png}{Despues de hacer clic al boton añadir se nos abre el modal.}
\captura{assets/img/ethoshub-profesional/04_experiencia/profesional_experiencia_laboral_5.png}{El primer campo nos permite registrar la empresa de la experiencia que queremos registrar.}
\captura{assets/img/ethoshub-profesional/04_experiencia/profesional_experiencia_laboral_6.png}{El segundo campo nos permite registrar el cargo que ocupamos en la experiencia laboral.}
\captura{assets/img/ethoshub-profesional/04_experiencia/profesional_experiencia_laboral_7.png}{El tercer campo nos permite registrar la fecha de inicio de la experiencia laboral.}
\captura{assets/img/ethoshub-profesional/04_experiencia/profesional_experiencia_laboral_8.png}{Tenemos 2 formas de registrar la fecha de inicio siendo el metodo manual o dandole click a la imagen de calendario.}
\captura{assets/img/ethoshub-profesional/04_experiencia/profesional_experiencia_laboral_9.png}{Al darle clicka la imagen de calendario nos aparece un mini-calendario donde podemos buscar la fecha que queremos registrar.}
\captura{assets/img/ethoshub-profesional/04_experiencia/profesional_experiencia_laboral_10.png}{El cuarto campo nos permite registrar la fecha de fin de la experiencia laboral.}
\captura{assets/img/ethoshub-profesional/04_experiencia/profesional_experiencia_laboral_11.png}{Tenemos 2 formas de registrar la fecha de fin siendo el metodo manual o dandole click a la imagen de calendario.}
\captura{assets/img/ethoshub-profesional/04_experiencia/profesional_experiencia_laboral_12.png}{Al darle clicka la imagen de calendario nos aparece un mini-calendario donde podemos buscar la fecha que queremos registrar.}
\captura{assets/img/ethoshub-profesional/04_experiencia/profesional_experiencia_laboral_13.png}{El quinto campo nos permite registrar la ubicacion de la empresa en la que se realizo la experiencia laboral.}
\captura{assets/img/ethoshub-profesional/04_experiencia/profesional_experiencia_laboral_14.png}{Tenemos 2 formas de registrar la ubicacion siendo el metodo manual insertando la direccion o dandole click a la imagen de mapa.}
\captura{assets/img/ethoshub-profesional/04_experiencia/profesional_experiencia_laboral_15.png}{Al hacer clic a la imagen de mapa nos aparece un modal donde se nos muestra el mapa.}
\captura{assets/img/ethoshub-profesional/04_experiencia/profesional_experiencia_laboral_16.png}{En el campo podemos ingresar manualmente la direccion de la empresa.}
\captura{assets/img/ethoshub-profesional/04_experiencia/profesional_experiencia_laboral_17.png}{El boton «cancelar» nos permite cancelar el registro de la direccion cerrando el modal del mapa.}
\captura{assets/img/ethoshub-profesional/04_experiencia/profesional_experiencia_laboral_18.png}{El boton «guardar» nos permite completar el registro del campo de dirrecion de la empresa cerrando el modal del mapa y volviendo al modal de registro.}
\captura{assets/img/ethoshub-profesional/04_experiencia/profesional_experiencia_laboral_19.png}{El sexto campo nos permite registrar el sitio web de la empresa de la que queremos registrar la experiencia laboral.}
\captura{assets/img/ethoshub-profesional/04_experiencia/profesional_experiencia_laboral_20.png}{El primer boton debajo del sexto campo «trabajo aqui actualmente» nos permite cambiar el registro , al seleccionarlo bloque el cuarto campo «fecha de fin» dejandolo vacio.}
\captura{assets/img/ethoshub-profesional/04_experiencia/profesional_experiencia_laboral_21.png}{El segundo boton debajo del sexto campo «Freelance/independiente» nos permite cambiar el registro, al seleccionarlo bloquea el primer campo «empresa» cambiando su contenido a trabajo independiente.}
\captura{assets/img/ethoshub-profesional/04_experiencia/profesional_experiencia_laboral_22.png}{El septimo campo nos permite registrar una descripcion de los logros alcanzados en la experiencia que queremos registrar.}
\captura{assets/img/ethoshub-profesional/04_experiencia/profesional_experiencia_laboral_23.png}{El octavo campo nos permite registrar las tecnologias que se usaron en la experiencia a registrar.}
\captura{assets/img/ethoshub-profesional/04_experiencia/profesional_experiencia_laboral_24.png}{El noveno campo (siendo un campo opcional) nos permite subir una imagen del logo de la empresaa la que pertenece la experiencia.}
\captura{assets/img/ethoshub-profesional/04_experiencia/profesional_experiencia_laboral_25.png}{Tenemos 2 metodos de registrar el logo de la empresa, el primero metodo nos permite subir el logo como archivo.}
\captura{assets/img/ethoshub-profesional/04_experiencia/profesional_experiencia_laboral_26.png}{Al hacer clic al boton subir archivo nos abre el explorador de archivos donde seleccionamos la imagen que se quiere subir.}
\captura{assets/img/ethoshub-profesional/04_experiencia/profesional_experiencia_laboral_27.png}{El segundo metodo nos permite subir el logo como un url.}
\captura{assets/img/ethoshub-profesional/04_experiencia/profesional_experiencia_laboral_28.png}{En el campo ingresamos la url de la imagen que queremos subir.}
\captura{assets/img/ethoshub-profesional/04_experiencia/profesional_experiencia_laboral_29.png}{Al darle al boton «aplicar» se guarda la url de la imagen mostrando una presualizacion de la misma.}
\captura{assets/img/ethoshub-profesional/04_experiencia/profesional_experiencia_laboral_30.png}{Tenemos 2 metodos de registrar el logo de la empresa, el primero metodo nos permite subir el banner como archivo.}
\captura{assets/img/ethoshub-profesional/04_experiencia/profesional_experiencia_laboral_31.png}{Al hacer clic al boton subir archivo nos abre el explorador de archivos donde seleccionamos la imagen que se quiere subir.}
\captura{assets/img/ethoshub-profesional/04_experiencia/profesional_experiencia_laboral_32.png}{El segundo metodo nos permite subir el banner como un url.}
\captura{assets/img/ethoshub-profesional/04_experiencia/profesional_experiencia_laboral_33.png}{En el campo ingresamos la url de la imagen que queremos subir.}
\captura{assets/img/ethoshub-profesional/04_experiencia/profesional_experiencia_laboral_34.png}{Al darle al boton «aplicar» se guarda la url de la imagen mostrando una presualizacion de la misma.}
\captura{assets/img/ethoshub-profesional/04_experiencia/profesional_experiencia_laboral_35.png}{El boton «cancelar» nos permite cancelar el registro de la experiencia laboral.}
\captura{assets/img/ethoshub-profesional/04_experiencia/profesional_experiencia_laboral_36.png}{El boton «guardar» nos permite completar y guardar el registro de la experiencia laboral.}
\captura{assets/img/ethoshub-profesional/04_experiencia/profesional_experiencia_laboral_37.png}{Una vez completado el registro se cierra el modal y se muestra la experiencia como se puede apreciar.}
\captura{assets/img/ethoshub-profesional/04_experiencia/profesional_experiencia_laboral_38.png}{Podemos apreciar como despues de que se completa el registro calcula y muestra el \# de años de experiencia, \# de empresas y \# de roles empleados.}
\captura{assets/img/ethoshub-profesional/04_experiencia/profesional_experiencia_laboral_39.png}{El boton con simbolo de ojo nos permite ver una previsualizacion de la experiencia registrada.}
\captura{assets/img/ethoshub-profesional/04_experiencia/profesional_experiencia_laboral_40.png}{Podemos apreciar como se ve la previsualizacion de la experiencia registrada.}
\captura{assets/img/ethoshub-profesional/04_experiencia/profesional_experiencia_laboral_41.png}{Los botones visibles nos permiten editar una experiencia registrada.}
\captura{assets/img/ethoshub-profesional/04_experiencia/profesional_experiencia_laboral_42.png}{El boton nos permite cerrar la previsualizacion de la experiencia registrada.}
\captura{assets/img/ethoshub-profesional/04_experiencia/profesional_experiencia_laboral_43.png}{El boton permite borrar una experiencia registrada.}
\captura{assets/img/ethoshub-profesional/04_experiencia/profesional_experiencia_laboral_44.png}{Podemos ver un boton que nos permite editar una experiencia registrada desde la pantalla principal.}
\captura{assets/img/ethoshub-profesional/04_experiencia/profesional_experiencia_laboral_45.png}{Podemos ver un boton que nos permite borrar una experiencia registrada desde la pantalla principal.}
\captura{assets/img/ethoshub-profesional/04_experiencia/profesional_experiencia_laboral_46.png}{Despues de darle a cualquiera de los botones de editar aparece el modal donde podemos editar todos los campos de la experiencia registrada.}
\captura{assets/img/ethoshub-profesional/04_experiencia/profesional_experiencia_laboral_47.png}{Despues de darle a cualquiera de los botones de editar aparece el modal donde podemos editar todos los campos de la experiencia registrada.}
\captura{assets/img/ethoshub-profesional/04_experiencia/profesional_experiencia_laboral_48.png}{El boton «cancelar» nos permite cancelar la edicion de los datos de una experiencia registrada.}
\captura{assets/img/ethoshub-profesional/04_experiencia/profesional_experiencia_laboral_49.png}{El boton «guardar» nos permite sobreescribir los cambios de una experiencia registrada.}
\captura{assets/img/ethoshub-profesional/04_experiencia/profesional_experiencia_laboral_50.png}{El boton «eliminar» nos permite borrar los cambios de una experiencia registrada.}
\captura{assets/img/ethoshub-profesional/04_experiencia/profesional_experiencia_laboral_51.png}{Despues de darle a cualquiera de los botones de «eliminar» nos devuelve a la pantalla de experiencia y nos aparece mensaje debajo de la experiencia registrada.}
\captura{assets/img/ethoshub-profesional/04_experiencia/profesional_experiencia_laboral_52.png}{Al darle al boton «eliminar» en el apartado de confirmacion esta se elimina de forma permanente desapareciendo del apartado deexperiencia del usurio.}
\captura{assets/img/ethoshub-profesional/04_experiencia/profesional_experiencia_laboral_53.png}{Al darle al boton «cancelar» en el apartado de confirmacion se cancela el intento de borrado de una experiencia, desapareciendo el apartado de confirmacion.}

\subsection{Registro de Formación Académica (Educación)}
\label{subsec:prof-educacion}
El bloque \campo{Educación} registra y muestra su formación y certificaciones.

\captura{assets/img/ethoshub-profesional/05_educacion/barra_lateral_profesional_educacion.png}{Para entrar al apartado, haga clic en \boton{Educación} (icono de gorra de graduado) en el panel lateral izquierdo.}
\captura{assets/img/ethoshub-profesional/05_educacion/Pagina_Principal_de_Educacion.png}{Vista principal del apartado \campo{Educación}, donde se concentra todo su historial académico.}
\captura{assets/img/ethoshub-profesional/05_educacion/Panel_Superior_Registra_Titulos.png}{El panel superior resume el número de títulos, los años de formación y la cantidad de centros registrados.}
\captura{assets/img/ethoshub-profesional/05_educacion/Panel_inferior_Historial_acedemico.png}{Para iniciar un registro pulse \boton{Historial académico} (arriba a la derecha) o el botón que aparece dentro del panel cuando aún no existe ningún registro.}
\captura{assets/img/ethoshub-profesional/05_educacion/Formulacio_registro_Nueva_trayectoria_vacio.png}{Al pulsar cualquiera de los botones se abre el formulario \campo{Nueva trayectoria}. Los campos marcados con asterisco (*) son obligatorios.}
\captura{assets/img/ethoshub-profesional/05_educacion/Formulacio_registro_Nueva_trayectoria_01.png}{En el campo \campo{Institución} escriba el nombre del centro donde se formó o estudia actualmente.}
\captura{assets/img/ethoshub-profesional/05_educacion/Formulacio_registro_Nueva_trayectoria_02.png}{El campo \campo{Tipo de educación} permite seleccionar la modalidad de estudios mediante un desplegable.}
\captura{assets/img/ethoshub-profesional/05_educacion/Formulacio_registro_Nueva_trayectoria_02a.png}{Al desplegar \campo{Tipo de educación} se muestran las opciones disponibles para clasificar la formación.}
\captura{assets/img/ethoshub-profesional/05_educacion/Educacion_combobox_tipo_educacion.png}{Seleccione en el desplegable el tipo de educación que corresponda a su trayectoria.}
\captura{assets/img/ethoshub-profesional/05_educacion/Formulacio_registro_Nueva_trayectoria_03.png}{En \campo{Título/Grado} indique el nombre del título o grado obtenido o en curso.}
\captura{assets/img/ethoshub-profesional/05_educacion/Formulacio_registro_Nueva_trayectoria_04.png}{En \campo{Área de estudio} especifique la especialidad o el área concreta de la formación.}
\captura{assets/img/ethoshub-profesional/05_educacion/Formulacio_registro_Nueva_trayectoria_05.png}{En \campo{Inicio y fin} registre las fechas de comienzo y conclusión de los estudios.}
\captura{assets/img/ethoshub-profesional/05_educacion/Formulacio_registro_Nueva_trayectoria_05a.png}{Al abrir el selector de fechas puede navegar por meses y años para elegir la fecha de inicio.}
\captura{assets/img/ethoshub-profesional/05_educacion/Formulacio_registro_Nueva_trayectoria_05b.png}{El calendario permite seleccionar el día exacto de inicio de la formación.}
\captura{assets/img/ethoshub-profesional/05_educacion/Formulacio_registro_Nueva_trayectoria_05c.png}{De igual forma, el selector de fecha de fin se despliega como un calendario interactivo.}
\captura{assets/img/ethoshub-profesional/05_educacion/Formulacio_registro_Nueva_trayectoria_05d.png}{Elija el mes y el año correspondientes a la fecha de finalización.}
\captura{assets/img/ethoshub-profesional/05_educacion/Formulacio_registro_Nueva_trayectoria_05e.png}{Confirme el día seleccionado para registrar la fecha de fin.}
\captura{assets/img/ethoshub-profesional/05_educacion/Formulacio_registro_Nueva_trayectoria_06.png}{En \campo{Promedio/GPA} introduzca la calificación obtenida, siguiendo el ejemplo que sugiere el sistema.}
\captura{assets/img/ethoshub-profesional/05_educacion/Formulacio_registro_Nueva_trayectoria_07.png}{En \campo{Link de verificación} pegue la dirección web del centro o de la credencial que acredita la formación.}
\captura{assets/img/ethoshub-profesional/05_educacion/Formulacio_registro_Nueva_trayectoria_08.png}{Active la casilla \campo{Estudio actualmente} si la formación sigue en curso.}
\captura{assets/img/ethoshub-profesional/05_educacion/Formulario_Bloqueo_Fin_fecha.png}{Al activar \campo{Estudio actualmente}, el campo de fecha de fin se bloquea automáticamente, ya que no procede indicarla.}
\captura{assets/img/ethoshub-profesional/05_educacion/Formulacio_registro_Nueva_trayectoria_09.png}{En \campo{Logo institución} puede subir el logotipo del centro desde un archivo local.}
\captura{assets/img/ethoshub-profesional/05_educacion/Formulacio_registro_Nueva_trayectoria_09a.png}{Al pulsar \boton{Subir archivo} se abre el explorador de archivos para seleccionar la imagen del logotipo.}
\captura{assets/img/ethoshub-profesional/05_educacion/Educacion_subir_logo_url.png}{Como alternativa, pulse \boton{URL} para cargar el logotipo a partir de una dirección web.}
\captura{assets/img/ethoshub-profesional/05_educacion/Educacion_subir_logo_url_01.png}{Pegue la URL del logotipo; el sistema mostrará una previsualización de la imagen.}
\captura{assets/img/ethoshub-profesional/05_educacion/Formulacio_registro_Nueva_trayectoria_10.png}{En \campo{Certificado/Título} puede subir el documento (imagen o PDF) que respalda la formación.}
\captura{assets/img/ethoshub-profesional/05_educacion/Formulacio_registro_Nueva_trayectoria_10a.png}{Al pulsar \boton{Subir archivo} se abre el explorador para seleccionar el certificado desde su equipo.}
\captura{assets/img/ethoshub-profesional/05_educacion/Educacion_subir_pdf_url.png}{También puede pulsar \boton{URL} para enlazar el certificado mediante una dirección web.}
\captura{assets/img/ethoshub-profesional/05_educacion/Educacion_subir_pdf_url_01.png}{Introduzca la URL del certificado para vincularlo al registro.}
\captura{assets/img/ethoshub-profesional/05_educacion/Formulacio_registro_Nueva_trayectoria_11.png}{Una vez completado el formulario, use \boton{Cancelar} para descartar o \boton{Guardar} para registrar la trayectoria.}
\captura{assets/img/ethoshub-profesional/05_educacion/Formulacio_registro_Nueva_trayectoria_11a.png}{El botón \boton{Guardar} confirma y almacena el nuevo registro de formación.}
\captura{assets/img/ethoshub-profesional/05_educacion/Formulacio_registro_Nueva_trayectoria_11b.png}{El botón \boton{Cancelar} cierra el formulario sin guardar los datos introducidos.}
\captura{assets/img/ethoshub-profesional/05_educacion/Educacion_panel_principal_registrado.png}{Tras guardar, la nueva trayectoria aparece reflejada en la página de Educación.}
\captura{assets/img/ethoshub-profesional/05_educacion/Educacion_Contador_de_registro.png}{El panel superior se actualiza: aumentan los contadores de títulos, años (calculados automáticamente a partir de las fechas) y centros.}
\captura{assets/img/ethoshub-profesional/05_educacion/Profesional_Educacion_registrado.png}{El historial académico muestra la trayectoria registrada en formato de listado.}
\captura{assets/img/ethoshub-profesional/05_educacion/Profesional_Educacion_registrado_01.png}{Cada registro presenta sus datos principales de forma resumida y ordenada.}
\captura{assets/img/ethoshub-profesional/05_educacion/Educacion_vista_previa_de_registro.png}{Al seleccionar un registro, a la derecha se despliega un panel con la previsualización de todos los campos completados.}
\captura{assets/img/ethoshub-profesional/05_educacion/Educacion_vista_previa_de_registro01.png}{Desde la vista previa puede pulsar el botón de edición para modificar cualquier campo del registro.}
\captura{assets/img/ethoshub-profesional/05_educacion/Educacion_vista_previa_de_registro02.png}{El enlace de verificación de credencial permite acceder al sitio web indicado en el formulario.}
\captura{assets/img/ethoshub-profesional/05_educacion/Educacion_vista_previa_de_registro03.png}{En \campo{Certificado adjunto} puede visualizar o descargar la imagen o el PDF cargado.}
\captura{assets/img/ethoshub-profesional/05_educacion/Educacion_vista_previa_de_registro04.png}{El botón de eliminación retira la certificación adjunta del registro.}
\captura{assets/img/ethoshub-profesional/05_educacion/Profesional_Educacion_registrado_02.png}{Desde la página principal también puede editar el historial académico con el botón de edición.}
\captura{assets/img/ethoshub-profesional/05_educacion/Profesional_Educacion_registrado_03.png}{El botón de papelera permite eliminar un registro del historial académico.}
\captura{assets/img/ethoshub-profesional/05_educacion/Educacion_eliminar_registro.png}{Antes de borrar, el sistema solicita confirmación con las opciones \boton{Cancelar} y \boton{Eliminar}.}

\subsection{Generación de Currículum con CV Studio}
\label{subsec:prof-cvstudio}
\campo{CV Studio} permite generar currículums profesionales a partir de plantillas en
formato LaTeX o Markdown, con apoyo de un asistente de inteligencia artificial, y
exportarlos o guardarlos en el perfil.

\captura{assets/img/ethoshub-profesional/06_cv-studio/CAP1.png}{Acceda a \boton{CV Studio} desde el panel lateral izquierdo para abrir el estudio de currículum, con sus botones, plantillas y editor.}
\captura{assets/img/ethoshub-profesional/06_cv-studio/CAP2.png}{El estudio ofrece dos formatos de trabajo: \boton{LaTeX} y \boton{Markdown}. En primer lugar seleccionamos \boton{LaTeX}.}
\captura{assets/img/ethoshub-profesional/06_cv-studio/CAP3.png}{Dentro de las plantillas LaTeX seleccionamos \campo{Moderno Pro}, que abre el editor con el contenido de la plantilla.}
\captura{assets/img/ethoshub-profesional/06_cv-studio/CAP4.png}{Vista ampliada del editor de la plantilla \campo{Moderno Pro}.}
\captura{assets/img/ethoshub-profesional/06_cv-studio/CAP5.png}{Seleccionamos ahora la plantilla LaTeX \campo{Ejecutivo Clásico}, que también dispone de su propio editor.}
\captura{assets/img/ethoshub-profesional/06_cv-studio/CAP6.png}{Vista ampliada del editor de la plantilla \campo{Ejecutivo Clásico}.}
\captura{assets/img/ethoshub-profesional/06_cv-studio/CAP7.png}{Seleccionamos la plantilla LaTeX \campo{Compacto Una Página} y se abre su editor correspondiente.}
\captura{assets/img/ethoshub-profesional/06_cv-studio/CAP8.png}{Vista ampliada del editor de la plantilla \campo{Compacto Una Página}.}
\captura{assets/img/ethoshub-profesional/06_cv-studio/CAP9.png}{Volvemos al selector de formato y ahora elegimos \boton{Markdown}.}
\captura{assets/img/ethoshub-profesional/06_cv-studio/CAP10.png}{Dentro de las plantillas Markdown seleccionamos \campo{Ingeniero Senior}; el editor muestra su contenido.}
\captura{assets/img/ethoshub-profesional/06_cv-studio/CAP11.png}{Vista ampliada del editor de \campo{Ingeniero Senior}; arriba a la derecha el zoom está al 100\,\%.}
\captura{assets/img/ethoshub-profesional/06_cv-studio/CAP12.png}{Seleccionamos la plantilla Markdown \campo{Stack \& Proyectos} y se abre su editor.}
\captura{assets/img/ethoshub-profesional/06_cv-studio/CAP13.png}{Vista ampliada del editor de \campo{Stack \& Proyectos}; el zoom está al 75\,\%.}
\captura{assets/img/ethoshub-profesional/06_cv-studio/CAP14.png}{Seleccionamos la plantilla Markdown \campo{Minimalista Senior} y se abre su editor.}
\captura{assets/img/ethoshub-profesional/06_cv-studio/CAP15.png}{Vista ampliada del editor de \campo{Minimalista Senior}; el zoom está al 100\,\%.}
\captura{assets/img/ethoshub-profesional/06_cv-studio/CAP16.png}{Al pulsar \boton{Nuevo CV} comenzamos un currículum desde cero en Markdown, con las mismas plantillas disponibles.}
\captura{assets/img/ethoshub-profesional/06_cv-studio/CAP17.png}{Al pulsar \boton{Nuevo CV} también podemos empezar desde cero en LaTeX, con las plantillas LaTeX disponibles.}
\captura{assets/img/ethoshub-profesional/06_cv-studio/CAP18.png}{El botón con icono de hojas pequeñas copia el código del editor al portapapeles.}
\captura{assets/img/ethoshub-profesional/06_cv-studio/CAP19.png}{Al pulsar \boton{Asistente IA} se abre una ventana de asistencia con inteligencia artificial.}
\captura{assets/img/ethoshub-profesional/06_cv-studio/CAP20.png}{En esta ventana podemos elegir entre las opciones sugeridas o escribir directamente la solicitud para enviarla a la IA.}
\captura{assets/img/ethoshub-profesional/06_cv-studio/CAP21.png}{Solicitamos a la IA un nuevo CV para un perfil de ingeniero en sistemas y la herramienta genera el código correspondiente.}
\captura{assets/img/ethoshub-profesional/06_cv-studio/CAP22.png}{El botón \boton{Aplicar al editor} carga automáticamente el resultado de la IA en nuestro editor.}
\captura{assets/img/ethoshub-profesional/06_cv-studio/CAP23.png}{Se muestra el nuevo CV generado por la IA; arriba a la derecha el zoom está al 75\,\%.}
\captura{assets/img/ethoshub-profesional/06_cv-studio/CAP24.png}{Al pulsar \boton{Exportar} podemos elegir el formato de exportación del currículum.}
\captura{assets/img/ethoshub-profesional/06_cv-studio/CAP25.png}{En el menú \boton{Exportar} seleccionamos \boton{Exportar PDF}.}
\captura{assets/img/ethoshub-profesional/06_cv-studio/CAP26.png}{Tras pulsar \boton{Exportar PDF}, el documento se genera y se descarga en formato PDF.}
\captura{assets/img/ethoshub-profesional/06_cv-studio/CAP27.png}{La opción \boton{Exportar Markdown} genera el archivo de la misma forma que la exportación a PDF.}
\captura{assets/img/ethoshub-profesional/06_cv-studio/CAP28.png}{La opción \boton{Exportar TXT} descarga el currículum en texto plano, igual que las exportaciones anteriores.}
\captura{assets/img/ethoshub-profesional/06_cv-studio/CAP29.png}{Al pulsar \boton{Guardar en mi perfil} se abre una ventana que solicita confirmar o cancelar la acción.}
\captura{assets/img/ethoshub-profesional/06_cv-studio/CAP30.png}{En la ventana \campo{Guardar documento} escribimos el nombre deseado y pulsamos \boton{Guardar} para almacenarlo en el perfil.}

\section{Mecanismos de Validación de Habilidades}
\label{sec:prof-habilidades}
El bloque \campo{Habilidades} permite declarar las competencias técnicas (hard skills)
y blandas (soft skills), asignarles un nivel y destacarlas para su validación pública
en la sección \campo{Top Skills}.

\subsection{Solicitud Formal de Validación Técnica (Hard Skills)}
\label{subsec:prof-hardskills}

\captura{assets/img/ethoshub-profesional/02_habilidades/profesional_habilidades_1.png}{Para ir al apartado de habilidades debe hacer clic en el botón de «Habilidades».}
\captura{assets/img/ethoshub-profesional/02_habilidades/profesional_habilidades_2.png}{Los botones «+ Hard Skill» y «+ Agregar primera habilidad» nos permite agregar una nueva hard skill.}
\captura{assets/img/ethoshub-profesional/02_habilidades/profesional_habilidades_3.png}{Despues de hacer clic en «+ Hard Skill» o + Agregar primera habilidad nos aparece un modal.}
\captura{assets/img/ethoshub-profesional/02_habilidades/profesional_habilidades_4.png}{El espacio de «Buscar tecnología...» nos permite realizar una búsqueda de una tecnología que tengamos como habilidad dura.}
\captura{assets/img/ethoshub-profesional/02_habilidades/profesional_habilidades_5.png}{El espacio de «Seleccionar nivel» nos permite elegir nuestro nivel.}
\captura{assets/img/ethoshub-profesional/02_habilidades/profesional_habilidades_6.png}{Los botones de «Cancelar» y «X» visibles nos permite cancelar y cerrar el modal para no registrar una habilidad.}
\captura{assets/img/ethoshub-profesional/02_habilidades/profesional_habilidades_7.png}{El boton de «Agregar» nos permite agregar una habilidad dura luego de completar los campos obligatorios.}

\subsection{Monitoreo y Seguimiento del Estado de la Solicitud}
\label{subsec:prof-topskills}
Tras registrar una habilidad técnica puede editar su nivel, destacarla en
\campo{Top Skills}, reordenar las destacadas y filtrarlas por categoría.

\captura{assets/img/ethoshub-profesional/02_habilidades/profesional_habilidades_8.png}{Una vez agregada una habilidad dura el boton con figura de un lapiz nos permite editar nuestra habilidad dura.}
\captura{assets/img/ethoshub-profesional/02_habilidades/profesional_habilidades_9.png}{En el modal podemos editar el nivel.}
\captura{assets/img/ethoshub-profesional/02_habilidades/profesional_habilidades_10.png}{Despues de editar nuestro nivel como la los botones de «Cancelar» y «X» nos permiten cancelar o cerrar la edicion.}
\captura{assets/img/ethoshub-profesional/02_habilidades/profesional_habilidades_11.png}{El boton de «Actualizar» nos permite actualizar nuestra habilidad dura.}
\captura{assets/img/ethoshub-profesional/02_habilidades/profesional_habilidades_12.png}{Despues de editar el boton de con figura de una estrella nos permite destacar nuestra habilidad dura.}
\captura{assets/img/ethoshub-profesional/02_habilidades/profesional_habilidades_13.png}{Despues de destacar nuestra habilidad dura se muestra en la seccion de «Top Skills».}
\captura{assets/img/ethoshub-profesional/02_habilidades/profesional_habilidades_14.png}{Despues de tener mas de una o dos habilidades duras destacadas los botones con flecha en direccion arriba y abajo nos permiten organizar por orden de prioridad de nuestra habilidades duras.}
\captura{assets/img/ethoshub-profesional/02_habilidades/profesional_habilidades_15.png}{Despues el boton de la figura con estrella nos permite quitar nuestra habilidad dura de la seccion de «Top Skills».}
\captura{assets/img/ethoshub-profesional/02_habilidades/profesional_habilidades_16.png}{Despues de apretar el boton de la estrella en la seccion de «Top Skills» la habilidad dura deja de estar en dicha seccion.}
\captura{assets/img/ethoshub-profesional/02_habilidades/profesional_habilidades_17.png}{Al darle a los botones «Todos, Frontend, Backend, Base de Datos, Infraestructura, Mobile, Diseño» visibles nos permiten filtrar las habilidades duras.}
\captura{assets/img/ethoshub-profesional/02_habilidades/profesional_habilidades_18.png}{Despues de darle al boton de «Frontend» se filtra las habilidades duras catalogadas como Frontend.}
\captura{assets/img/ethoshub-profesional/02_habilidades/profesional_habilidades_19.png}{Luego al darle a otro boton en este caso «Base de Datos» se filtra las habilidades duras catalogadas como Base de Datos.}

\subsection{Registro de Habilidades Blandas (Soft Skills)}
\label{subsec:prof-softskills}

\captura{assets/img/ethoshub-profesional/02_habilidades/profesional_habilidades_20.png}{Despues los botones de «+ Soft Skill» y «+ Agregar» visibles nos permiten agregar una nueva habilidad blanda.}
\captura{assets/img/ethoshub-profesional/02_habilidades/profesional_habilidades_21.png}{Despues de darle aparece un modal donde en la seccion de «Titulo» podemos ingresar un titulo de nuestra habilidad blanda.}
\captura{assets/img/ethoshub-profesional/02_habilidades/profesional_habilidades_22.png}{En el espacio de «Caso de éxito» podemos ingresar un escenario donde tuvimos éxito aplicando la habilidad blanda con un maximo de 250 caracteres.}
\captura{assets/img/ethoshub-profesional/02_habilidades/profesional_habilidades_23.png}{Los botones de «Cancelar» y «X» visibles nos permiten cancelar y cerrar el modal.}
\captura{assets/img/ethoshub-profesional/02_habilidades/profesional_habilidades_24.png}{El boton de «Agregar» nos permite agregar una nueva habilidad blanda luego de completar los campos obligatorios.}
\captura{assets/img/ethoshub-profesional/02_habilidades/profesional_habilidades_25.png}{Despues de agregar una habilidad blanda es visible.}
\captura{assets/img/ethoshub-profesional/02_habilidades/profesional_habilidades_26.png}{El boton con la figura de un lapiz nos permite editar nuestra habilidad blanda.}
\captura{assets/img/ethoshub-profesional/02_habilidades/profesional_habilidades_27.png}{Despues de hacer clic en el boton con figura de lapiz podemos editar los campos de «Titulo» y «Caso de éxito».}
\captura{assets/img/ethoshub-profesional/02_habilidades/profesional_habilidades_28.png}{Los boton de «Cancelar» y «X» nos permiten cancelar o cerrar el modal de «Editar Soft Skill» como en la.}
\captura{assets/img/ethoshub-profesional/02_habilidades/profesional_habilidades_29.png}{El boton de «Actualizar» nos permite actualizar nuestra habilidad blanda una vez que terminemos de editar los campos obligatorios.}
\captura{assets/img/ethoshub-profesional/02_habilidades/profesional_habilidades_30.png}{El boton con icono de un tacho de basura nos permite eliminar una habilidad blanda.}
\captura{assets/img/ethoshub-profesional/02_habilidades/profesional_habilidades_31.png}{El boton con icono de un tacho de basura nos permite eliminar una habilidad dura.}

\section{Configuración y Privacidad del Perfil Público}
\label{sec:prof-configuracion}
Desde \campo{Configuración} se gestionan la identidad del perfil, la personalización de
la interfaz, la seguridad de la cuenta y la exportación o eliminación de datos.

\subsection{Identidad del Perfil}
\label{subsec:prof-identidad}

\captura{assets/img/ethoshub-profesional/09_configuracion/configuracion_profesional_1.png}{Para ir al apartado de configuración debe hacer clic en el botón de «Configuración».}
\captura{assets/img/ethoshub-profesional/09_configuracion/configuracion_profesional_2.png}{El botón de «Identidad» muestra la sección entera para configurar el perfil.}
\captura{assets/img/ethoshub-profesional/09_configuracion/configuracion_profesional_3.png}{Las secciones de «Foto de perfil» se remarcan en 3 lugares distintos.}
\captura{assets/img/ethoshub-profesional/09_configuracion/configuracion_profesional_4.png}{El boton de «Seleccionar Imagen» y el campo de «pega una URL de imagen» nos permite cambiar nuestra foto de perfil.}
\captura{assets/img/ethoshub-profesional/09_configuracion/configuracion_profesional_5.png}{El botón de «Guardar» nos permite guardar nuestra foto de perfil.}
\captura{assets/img/ethoshub-profesional/09_configuracion/configuracion_profesional_6.png}{Al hacer clic en «Guardar» vemos que las 3 secciones donde se visualiza la foto de perfil se que el cambio.}
\captura{assets/img/ethoshub-profesional/09_configuracion/configuracion_profesional_7.png}{Los campos de «Nombre» y «Apellido» nos permiten colocar nuestro nombres y apellidos respectivo con un valor máximo de 80 caracteres.}
\captura{assets/img/ethoshub-profesional/09_configuracion/configuracion_profesional_8.png}{En el campo de «Nivel profesional» podemos seleccionar el nivel que tenemos al hacer clic en cualquiera de las opciones.}
\captura{assets/img/ethoshub-profesional/09_configuracion/configuracion_profesional_9.png}{En el campo de «Disponibilidad» podemos seleccionar al hacer clic en cualquiera de las opciones.}
\captura{assets/img/ethoshub-profesional/09_configuracion/configuracion_profesional_10.png}{En el campo de «Ubicación» podemos escribir nuestra dirección actual con un valor máximo de 80 caracteres.}
\captura{assets/img/ethoshub-profesional/09_configuracion/configuracion_profesional_11.png}{El botón con icono de un marcador nos permite desplegar un mapa interactivo.}
\captura{assets/img/ethoshub-profesional/09_configuracion/configuracion_profesional_12.png}{Luego de hacer click en el icono de marcador vemos un modal de «Seleccionar ubicación».}
\captura{assets/img/ethoshub-profesional/09_configuracion/configuracion_profesional_13.png}{En el campo de «Busca una dirección o lugar...» podemos buscar un lugar mediante calles, avenida, ciudad o país.}
\captura{assets/img/ethoshub-profesional/09_configuracion/configuracion_profesional_14.png}{En el mapa interactivo si hacemos clic en alguna sección del mapa se muestra un icono de marcador que podemos arrastrar para marcar nuestra ubicación exacta.}
\captura{assets/img/ethoshub-profesional/09_configuracion/configuracion_profesional_15.png}{El botón de «Cancelar» o «X» que es nos permite al hacer clic cancelar o cerrar el modal de selección de ubicación.}
\captura{assets/img/ethoshub-profesional/09_configuracion/configuracion_profesional_16.png}{El botón de «Confirmar» nos permite confirmar nuestra ubicación y guardar con un autocompletado automatico.}
\captura{assets/img/ethoshub-profesional/09_configuracion/configuracion_profesional_17.png}{El campo de «Sitio web personal» nos permite ingresar un sitio web personal que tengamos con un máximo de 80 caracteres.}
\captura{assets/img/ethoshub-profesional/09_configuracion/configuracion_profesional_18.png}{El botón «Guardar perfil» al hacer clic guardar nuestros datos profesionales.}
\captura{assets/img/ethoshub-profesional/09_configuracion/configuracion_profesional_19.png}{Vemos el llenado de nuestros datos profesionales.}
\captura{assets/img/ethoshub-profesional/09_configuracion/configuracion_profesional_20.png}{Los datos de profesionales son posibles actualizarlos una vez que modifiques los campos y le demos clic al botón «Guardar perfil».}
\captura{assets/img/ethoshub-profesional/09_configuracion/configuracion_profesional_21.png}{En la sección de «Biografia profesional» podemos ingresar una pequeña biografia con un máximo de 500 caracteres.}
\captura{assets/img/ethoshub-profesional/09_configuracion/configuracion_profesional_22.png}{El botón de «Guardar bio» nos permite guardar nuestra biografia personal.}
\captura{assets/img/ethoshub-profesional/09_configuracion/configuracion_profesional_23.png}{Una vez completados todos los campos podemos observar que se guardar correctamente.}

\subsection{Personalización de la Interfaz}
\label{subsec:prof-personalizacion}

\captura{assets/img/ethoshub-profesional/09_configuracion/configuracion_profesional_24.png}{Para ir a la sección de personalización debemos hacer clic en «Personalización».}
\captura{assets/img/ethoshub-profesional/09_configuracion/configuracion_profesional_25.png}{En la sección de «Tema de la Interfaz» podemos hacer clic en «Claro» para seleccionar un modo claro de la plataforma.}
\captura{assets/img/ethoshub-profesional/09_configuracion/configuracion_profesional_26.png}{El botón de «Oscuro» al hacer clic cambia a modo oscuro toda la plataforma y el botón «Sistema» se adapta al tema del Navegador.}
\captura{assets/img/ethoshub-profesional/09_configuracion/configuracion_profesional_27.png}{En la sección de «Idioma de la Interfaz» podemos hacer clic en «ES» «EN» «BR» para seleccionar un idioma a nuestra elección.}

\subsection{Seguridad de la Cuenta}
\label{subsec:prof-seguridad}

\captura{assets/img/ethoshub-profesional/09_configuracion/configuracion_profesional_28.png}{Para ir a seguridad debemos hacer clic en «Seguridad».}
\captura{assets/img/ethoshub-profesional/09_configuracion/configuracion_profesional_29.png}{En la sección de «Cambiar contraseña» en el botón de «Solicitar cambio de contraseña» podemos hacer clic para pedir un código en el correo asociado al momento de registrarse en la plataforma.}
\captura{assets/img/ethoshub-profesional/09_configuracion/configuracion_profesional_30.png}{Luego de hacer clic en la sección de «Codigo de verificación» debemos ingresar un código que nos mando la plataforma a nuestro correo asociado.}
\captura{assets/img/ethoshub-profesional/09_configuracion/configuracion_profesional_31.png}{El botón de «Volver» nos permite volver y cancelar la accion de cambiar contraseña.}
\captura{assets/img/ethoshub-profesional/09_configuracion/configuracion_profesional_32.png}{El botón «Verificar código» nos permite verificar el código ingresado para avanzar al siguiente paso.}
\captura{assets/img/ethoshub-profesional/09_configuracion/configuracion_profesional_33.png}{Podemos ver el código que la plataforma envió a nuestro correo asociado a la cuenta.}
\captura{assets/img/ethoshub-profesional/09_configuracion/configuracion_profesional_34.png}{Luego de avanzar al siguiente paso tenemos 2 campos donde debemos ingresar la nueva contraseña.}
\captura{assets/img/ethoshub-profesional/09_configuracion/configuracion_profesional_35.png}{Los botones con icono de ojos visibles al hacer clic nos permiten ver la contraseña que ingresamos y Tambien ocultar.}
\captura{assets/img/ethoshub-profesional/09_configuracion/configuracion_profesional_36.png}{El botón de «Volver» nos permite volver y cancelar la accion de cambiar contraseña.}
\captura{assets/img/ethoshub-profesional/09_configuracion/configuracion_profesional_37.png}{El botón «Actualizar contraseña» nos permite cambiar nuestra contraseña por la nueva ingresada.}
\captura{assets/img/ethoshub-profesional/09_configuracion/configuracion_profesional_38.png}{En la sección de «Cambiar correo electrónico» en el botón de «Solicitar cambio de correo» al hacer clic nos envia un código a nuestro correo asociado.}
\captura{assets/img/ethoshub-profesional/09_configuracion/configuracion_profesional_39.png}{En el campo de «Codigo de verificación» debemos ingresar el código que nos mando la plataforma al correo.}
\captura{assets/img/ethoshub-profesional/09_configuracion/configuracion_profesional_40.png}{Se ve el código que la plataforma nos envió para el cambio de correo electrónico.}
\captura{assets/img/ethoshub-profesional/09_configuracion/configuracion_profesional_41.png}{El botón de «Volver» nos permite volver y cancelar la accion de cambiar correo electrónico.}
\captura{assets/img/ethoshub-profesional/09_configuracion/configuracion_profesional_42.png}{El botón «Verificar código» nos permite verificar el código ingresado para avanzar al siguiente paso.}
\captura{assets/img/ethoshub-profesional/09_configuracion/configuracion_profesional_43.png}{Luego de avanzar al siguiente paso tenemos un campo donde debemos ingresar un nuevo correo electrónico.}
\captura{assets/img/ethoshub-profesional/09_configuracion/configuracion_profesional_44.png}{El botón de «Volver» nos permite volver y cancelar la accion de cambiar correo electrónico.}
\captura{assets/img/ethoshub-profesional/09_configuracion/configuracion_profesional_45.png}{El botón «Actualizar correo» al hacer clic nos permite cambiar nuestro correo por el nuevo ingresado.}

\subsection{Exportación de Datos y Zona de Peligro}
\label{subsec:prof-datos}

\captura{assets/img/ethoshub-profesional/09_configuracion/configuracion_profesional_46.png}{En la sección de «Exportar mis datos» el botón visible «Exportar datos» al hacer clic permite que podamos descargar nuestros datos completos.}
\captura{assets/img/ethoshub-profesional/09_configuracion/configuracion_profesional_47.png}{Se visualiza un modal de cargar de la exportación de datos.}
\captura{assets/img/ethoshub-profesional/09_configuracion/configuracion_profesional_48.png}{Vista del apartado dentro de la plataforma.}
\captura{assets/img/ethoshub-profesional/09_configuracion/configuracion_profesional_49.png}{Una vez que termine de cargar la exportación de datos nos permite elegir la ubicación donde guardaremos nuestros datos con el nombre de la plataforma y fecha.}
\captura{assets/img/ethoshub-profesional/09_configuracion/configuracion_profesional_50.png}{Luego de que aparezca el modal los botones «Cancelar» y «Enviar código» visibles que al hacer clic en cancelar cancela y cierra el modal y al hacer clic en enviar código nos manda al correo asociado.}
\captura{assets/img/ethoshub-profesional/09_configuracion/configuracion_profesional_51.png}{Se visualiza el código que la plataforma nos mando al correo asociado.}
\captura{assets/img/ethoshub-profesional/09_configuracion/configuracion_profesional_52.png}{Los botones «Volver» y «Confirmar eliminación» una vez agregado el código nos permite volver la accion de eliminar cuenta y el de confirmar eliminación para eliminar la cuenta.}

\begin{AvisoAdvertencia}
La eliminación de la cuenta es irreversible y borra todo su portafolio. El sistema exige
un código de verificación enviado a su correo antes de confirmar la acción.
\end{AvisoAdvertencia}

\section{Acceso desde Dispositivos Móviles (Diseño Adaptable)}
\label{sec:prof-responsive}
La plataforma cuenta con diseño adaptable (\textit{responsive}): todas las secciones del
perfil Profesional se reorganizan automáticamente para su uso cómodo en teléfonos y
tabletas. A continuación se muestran las vistas móviles más representativas.

\responsivePar{assets/img/ethoshub-profesional/10_responsive/profesional_responsive_mi_portafolio.jpg}{Mi portafolio en vista móvil.}{assets/img/ethoshub-profesional/10_responsive/profesional_responsive_habilidades.jpg}{Apartado de Habilidades en vista móvil.}
\responsivePar{assets/img/ethoshub-profesional/10_responsive/profesional_responsive_habilidades_hard.jpg}{Alta de una habilidad técnica (hard skill) en móvil.}{assets/img/ethoshub-profesional/10_responsive/profesional_responsive_habilidades_soft.jpg}{Alta de una habilidad blanda (soft skill) en móvil.}
\responsivePar{assets/img/ethoshub-profesional/10_responsive/profesional_responsive_proyectos.jpg}{Apartado de Proyectos en vista móvil.}{assets/img/ethoshub-profesional/10_responsive/profesional_responsive_formulario_proyectos.jpg}{Formulario de creación de proyecto en móvil.}
\responsivePar{assets/img/ethoshub-profesional/10_responsive/profesional_responsive_experiencia.jpg}{Apartado de Experiencia en vista móvil.}{assets/img/ethoshub-profesional/10_responsive/profesional_responsive_formulario_experiencia.jpg}{Formulario de experiencia en móvil (parte 1).}
\responsivePar{assets/img/ethoshub-profesional/10_responsive/profesional_responsive_formulario_experiencia2.jpg}{Formulario de experiencia en móvil (parte 2).}{assets/img/ethoshub-profesional/10_responsive/profesional_responsive_educacion.jpg}{Apartado de Educación en vista móvil.}
\responsivePar{assets/img/ethoshub-profesional/10_responsive/profesional_responsive_formulario_educacion.jpg}{Formulario de educación en móvil (parte 1).}{assets/img/ethoshub-profesional/10_responsive/profesional_responsive_formulario_educacion2.jpg}{Formulario de educación en móvil (parte 2).}
\responsivePar{assets/img/ethoshub-profesional/10_responsive/profesional_responsive_CVstudio.jpg}{CV Studio en vista móvil.}{assets/img/ethoshub-profesional/10_responsive/profesional_responsive_chat.jpg}{Bandeja de chat en vista móvil.}
\responsivePar{assets/img/ethoshub-profesional/10_responsive/profesional_responsive_chat_iniciado.jpg}{Conversación de chat iniciada en móvil.}{assets/img/ethoshub-profesional/10_responsive/profesional_responsive_conexiones.jpg}{Apartado de Conexiones en vista móvil.}
\responsivePar{assets/img/ethoshub-profesional/10_responsive/profesional_responsive_configuracion.jpg}{Configuración en vista móvil.}{assets/img/ethoshub-profesional/10_responsive/profesional_responsive_configuracion_campos.jpg}{Edición de campos de configuración en móvil.}
\responsivePar{assets/img/ethoshub-profesional/10_responsive/profesional_responsive_configuracion_personalizacion.jpg}{Pestaña de Personalización en móvil.}{assets/img/ethoshub-profesional/10_responsive/profesional_responsive_configuracion_seguridad.jpg}{Pestaña de Seguridad en móvil.}
\responsiveUna{assets/img/ethoshub-profesional/10_responsive/profesional_responsive_barra_lateral.jpg}{Menú lateral desplegable en vista móvil.}
